%======================================================================
%  Companion to Chapter 6 -- From prediction-correction to LATIN and
%  the direct cyclic method
%======================================================================
\chapter{From prediction--correction to LATIN and the direct cyclic method}\label{cc:ch6}

\framepart{Outline}

Chapter~\ref{B-ch:pl-cyc} of the notes works on two structures made of
fibres, the three-bar truss and the Bree vessel, where every algorithm can
be written in a few lines. Cast3M brings two things: the same fibre structures
with its own return maps, and the step to a real geometry, where the regimes of
the fibre model are only approximate. It also contains two of the methods of
the chapter as procedures, \texttt{@STATIO} and \texttt{@ZACPLUS}. This chapter
computes Bree's problem on an axisymmetric tube cycle by cycle, and lists what
remains to be done with the fields of Cast3M.

\section{Where Cast3M enters Chapter 6}\label{cc:sec6-map}
\index{Cast3M!Chapter 6}%

\begin{gbox}{Chapter 6 of the notes and Cast3M}\label{cc:tab6}
\small
\begin{tabular}{@{}p{42mm}@{\hspace{3mm}}p{98mm}@{}}
Bree vessel, Sec.~\ref{B-sec:cy-models}; Ex.~\ref{B-exo:cy-bree}
 & \textbf{Exercise \ref{exo:cc-bree}}: a thin tube under pressure and a cyclic gradient of
   temperature, \op{pasapas} cycle by cycle; ratchet per cycle against the
   fibre model.\\[1mm]
Three-bar truss, regimes, Ex.~\ref{B-exo:cy-truss}
 & With \texttt{cast3m/ch7/truss.template} and the loadings (1)--(3) of the
   chapter: shakedown, alternating plasticity, ratchet $(a,-a,0)$,
   $a=0.4\delta_Y$.\\[1mm]
Initial-strain iteration, Box~6.1
 & It is the equilibrium loop of Exercise~\ref{B-exo:pn-fields}: elastic stiffness assembled
   once, return map from the converged state.\\[1mm]
Whole history, Box~6.2; direct cyclic method, Sec.~\ref{B-sec:cy-dcm}
 & To do: the plastic strain history stored as a table of Gauss-point fields
   $\epsp(t_n)$; global stage: one factorization, $N$ elastic solves with the
   initial strain; local stage: \texttt{rrj2} swept along the history;
   closure $\vz(t_0)\leftarrow\vz(t_N)$.\\[1mm]
Period map, Sec.~\ref{B-sec:cy-period}
 & To do: one cycle of \op{pasapas} restarted from a saved state is the map
   $\Pi$; Krasnosel'skii--Mann averaging written on the fields; Anderson
   acceleration by an outside driver, as in Chapter~\ref{cc:ch7}.\\[1mm]
Zarka, Box~6.5; stationary methods, Box~6.4
 & To do: \texttt{@ZACPLUS} on the Bree tube at $(X,Y)=(0.6,1.5)$ and
   $(0.7,1.5)$, $H=0.02E$, against the cycle-by-cycle limit;
   \texttt{@STATIO} on its rail example, as an illustration of Box~6.4.\\
\end{tabular}
\end{gbox}

\section{A cyclic thermal load in \texttt{pasapas}}\label{cc:sec6-thermal}
\index{Bree diagram!Cast3M}%
\begin{castemcom}
The temperature is a nodal field \texttt{T} linear across the wall,
$\tau=-\tau_0x$ with $x=(r-r_m)/e$, times the history $\lambda(t)=\abs{\sin t}$.
\op{char} \texttt{'T'} makes it a thermal loading: \op{pasapas} computes the
thermal strain $\alpha\tau$ with the coefficient \texttt{ALPH} of the material.
The pressure is a mechanical loading constant in time. The top face stays
plane and free of axial force: one common \texttt{UZ}, \op{rela}
\texttt{ense}.
\end{castemcom}
\begin{castemlines}
\begin{verbatim}
rr  = coor 1 dom;
un  = rr masq 'SUPERIEUR' 0.;
tau = exco ((-1. * (tau0 / e))
      * (rr - (rm * un))) 'SCAL' 'T';
llam = abs (sin ((180. / pi0) * ltt));
cht = char 'T' tau
      (evol manu 'TEMPS' ltt 'LAMB' llam);
chp = char 'MECA' fp
      (evol manu 'TEMPS' ltt 'P' lun);
aa = (bloq uz bas) et (rela ense uz haut);
\end{verbatim}
\end{castemlines}

\framepart{Summary}

\begin{gbox*}
\begin{itemize}
\item \textbf{A cyclic computation is a long \op{pasapas}}: the regime is
  read from the drift of the mean strain cycle after cycle.
\item \textbf{The fibre model is a limit}: on a real wall the thermal stress is
  biaxial and the regime boundaries of the Bree diagram move; the tube shows by
  how much.
\item \textbf{The direct methods of the chapter act on histories of fields}:
  tables of Gauss-point fields indexed by the instants are the data structure
  of Box~6.2 and of the period map in gibiane.
\end{itemize}
\end{gbox*}

\framepart{Exercises}

\exercise{Bree's problem on a tube}\label{exo:cc-bree}
\index{Bree vessel!Cast3M}%
A thin tube ($r_m=100$\,mm, $e=5$\,mm, $E=200$\,GPa, $\nu=0.3$,
$\sY=280$\,MPa, $\alpha=10^{-5}$\,K$^{-1}$) is open, free to stretch, under a
constant internal pressure and a temperature linear across the wall,
$\tau=-\tau_0x\abs{\sin t}$. With $X=\sigma_P/\sY$, $\sigma_P=pr_m/e$, and
$Y=\sigma_T/\sY$, $\sigma_T=E\alpha\tau_0/(2(1-\nu))$, compute ten cycles for
$(X,Y)=(0.7,1.5)$, $(0.6,1.5)$ and $(0.2,3)$ with perfect plasticity, and
read the regime from the mean hoop strain at the end of each cycle. Compare
with the fibre model of Exercise~\ref{B-exo:cy-bree}.
\begin{solution}
The fibre model ($\tilde E=E/(1-\nu)$, $H=0$) ratchets at $(0.7,1.5)$ by
$0.3167\varepsilon_Y$ per cycle, shakes down at $(0.6,1.5)$ and alternates at
$(0.2,3)$ (\texttt{exo\_bree\_cyclic.py}). In the tube the thermal stress is
biaxial, as in the fibre model with $\tilde E=E/(1-\nu)$, but the pressure
stress is hoop only: the von Mises criterion then sees a different
combination, and the regime boundaries move. The run decides whether
$(0.6,1.5)$, close to the boundary $X+Y/4=1$ of the diagram, still shakes
down; the ratchet at $(0.7,1.5)$ should be of the same order as in the fibre
model, constant from the second cycle on, and the residual stress at the end
of a cycle should no longer change.
\untested{\texttt{cast3m/ch6/bree\_tube.dgibi}}
\lstinputlisting[style=gibstyle,firstline=29]{cast3m/ch6/bree_tube.dgibi}
\end{solution}

\begin{chapterbib}{9}
\bibitem{cc-bree} J.~Bree, Elastic--plastic behaviour of thin tubes
  subjected to internal pressure and intermittent high-heat fluxes with
  application to fast-nuclear-reactor fuel elements, \emph{Journal of Strain
  Analysis} 2 (1967) 226--238.
\bibitem{cc-zacplus} Cast3M, notices of the procedures \texttt{@ZACPLUS}
  and \texttt{@STATIO}, \url{https://www-cast3m.cea.fr}.
\end{chapterbib}
